% ============================================================
% CHAPITRE 25 — Branding et systèmes d'identité visuelle
% ============================================================

\chapter{Branding et systèmes d'identité visuelle}

\noindent
\textit{Dans les chapitres précédents, nous avons appris à créer des affiches, des flyers, des visuels pour les réseaux sociaux, et à conduire des campagnes visuelles. Tous ces supports reposent sur un socle commun : l'\emph{identité visuelle} d'une marque. Nous avons déjà abordé l'identité visuelle dans la troisième partie, à travers la création de logos et d'éléments graphiques. Mais il est temps d'aller plus loin et d'explorer la discipline qui englobe tout cela : le \emph{branding}.}

\vspace{1em}

\noindent
Le \terme{branding} est bien plus qu'un logo, plus qu'une palette de couleurs, plus qu'une typographie : c'est l'\emph{ensemble des perceptions, des émotions et des associations} qu'une marque suscite chez son public. Le branding ne se limite pas au design : il englobe la stratégie, le positionnement, la communication, l'expérience client. Mais le design graphique en est l'un des piliers : c'est lui qui donne une forme visible, tangible, mémorable à la marque.

\vspace{1em}

\noindent
Un \terme{système d'identité visuelle} est l'ensemble cohérent des éléments graphiques qui permettent de reconnaître une marque : logo, couleurs, typographies, éléments secondaires, règles d'application. Ce système doit être pensé, construit, documenté, pour être utilisé de manière cohérente sur tous les supports. Une identité visuelle réussie est reconnaissable, cohérente, intemporelle, adaptable. Une identité visuelle ratée est incohérente, générique, datée, difficile à décliner.

\vspace{1em}

\noindent
Ce chapitre a pour objectif de te faire maîtriser le branding et les systèmes d'identité visuelle. Nous commencerons par définir la marque et l'identité visuelle. Nous distinguerons le logo, l'identité visuelle et le branding. Nous apprendrons à effectuer une recherche et un positionnement visuel, à créer un logo et ses variantes, à choisir les couleurs de marque, à définir le système typographique, à créer des éléments graphiques secondaires, à appliquer l'identité sur les supports, à élaborer une charte graphique, et à présenter un projet d'identité visuelle. À la fin de ce chapitre, tu sauras construire des marques fortes et cohérentes, de la stratégie à la livraison.

% ============================================================
\section{Définition de la marque et de l'identité visuelle}
% ============================================================

\subsection{Qu'est-ce qu'une marque ?}

\begin{quote}
\textbf{Une marque est un ensemble de signes, de valeurs, de promesses et d'émotions qui permettent d'identifier un produit, un service, une organisation, et de le distinguer de ses concurrents.}
\end{quote}

Une marque n'est pas seulement un nom ou un logo : c'est une \terme{construction mentale}. Elle existe dans l'esprit des consommateurs, à travers :

\begin{itemize}
    \item \terme{Son} nom (Coca-Cola, Apple, Nike).
    \item \terme{Son} logo (pomme croquée, swoosh).
    \item \terme{Ses} couleurs (rouge, bleu, vert).
    \item \terme{Ses} typographies (cursive, sans-serif).
    \item \terme{Ses} valeurs (liberté, innovation, tradition).
    \item \terme{Sa} personnalité (jeune, sérieuse, élégante).
    \item \terme{Ses} promesses (plaisir, performance, durabilité).
    \item \terme{Son} expérience (service, qualité, relation).
\end{itemize}

La marque est ce qui reste quand on a tout oublié : le souvenir, l'émotion, l'association.

\subsection{Les composantes d'une marque}

\subsubsection{Le nom}

Le \terme{nom} est l'élément verbal de la marque. Il doit être :

\begin{itemize}
    \item \terme{Mémorable} : facile à retenir.
    \item \terme{Distinctif} : unique, différenciant.
    \item \terme{Prononçable} : facile à prononcer.
    \item \terme{Disponible} : libre juridiquement.
    \item \terme{Pertinent} : en lien avec l'activité.
\end{itemize}

\subsubsection{Le logo}

Le \terme{logo} est l'élément visuel de la marque. Il doit être :

\begin{itemize}
    \item \terme{Simple} : peu de formes, peu de couleurs.
    \item \terme{Mémorable} : reconnaissable, unique.
    \item \terme{Intemporel} : ne se démode pas.
    \item \terme{Versatile} : fonctionne sur tous les supports.
    \item \terme{Pertinent} : en lien avec les valeurs de la marque.
\end{itemize}

\subsubsection{Les couleurs}

Les \terme{couleurs} sont un élément fort de la marque. Elles doivent être :

\begin{itemize}
    \item \terme{Cohérentes} : avec les valeurs, le positionnement.
    \item \terme{Distinctives} : différentes des concurrents.
    \item \terme{Mémorables} : associées à la marque.
    \item \terme{Adaptées} : au public, à la culture.
\end{itemize}

\subsubsection{Les typographies}

Les \terme{typographies} sont un élément de personnalité de la marque. Elles doivent être :

\begin{itemize}
    \item \terme{Lisibles} : à toutes les tailles.
    \item \terme{Cohérentes} : avec le ton, l'ambiance.
    \item \terme{Distinctives} : reconnaissables.
    \item \terme{Versatiles} : utilisables sur tous les supports.
\end{itemize}

\subsubsection{Les valeurs}

Les \terme{valeurs} sont les principes qui guident la marque. Elles définissent son comportement, ses choix, sa communication.

\begin{itemize}
    \item \terme{Exemples} : innovation, tradition, écologie, qualité, accessibilité.
\end{itemize}

\subsubsection{La personnalité}

La \terme{personnalité} est le caractère de la marque. Elle définit son ton, son style, sa manière d'être.

\begin{itemize}
    \item \terme{Exemples} : jeune, sérieuse, élégante, rebelle, chaleureuse.
\end{itemize}

\subsubsection{La promesse}

La \terme{promesse} est ce que la marque s'engage à offrir à son public. Elle doit être claire, crédible, différenciante.

\begin{itemize}
    \item \terme{Exemples} : plaisir, performance, durabilité, simplicité.
\end{itemize}

\subsection{Qu'est-ce que l'identité visuelle ?}

\begin{quote}
\textbf{L'identité visuelle est l'ensemble des éléments graphiques qui permettent d'identifier et de représenter une marque : logo, couleurs, typographies, éléments secondaires, règles d'application.}
\end{quote}

L'identité visuelle est la \terme{traduction graphique} de la marque. Elle rend visible, tangible, mémorable ce qui est abstrait : les valeurs, la personnalité, la promesse.

Elle comprend :

\begin{itemize}
    \item \terme{Le} logo et ses variantes.
    \item \terme{La} palette de couleurs.
    \item \terme{Les} typographies.
    \item \terme{Les} éléments graphiques secondaires (icônes, pictogrammes, motifs).
    \item \terme{Les} règles d'application (charte graphique).
    \item \terme{Les} déclinaisons sur les supports (papeterie, web, packaging, signalétique).
\end{itemize}

\subsection{Les qualités d'une bonne identité visuelle}

\subsubsection{La reconnaissance}

Une bonne identité visuelle est \terme{reconnaissable} : on identifie la marque immédiatement, sans effort.

\subsubsection{La cohérence}

Une bonne identité visuelle est \terme{cohérente} : tous les éléments s'accordent, forment un tout.

\subsubsection{La distinction}

Une bonne identité visuelle est \terme{distinctive} : elle se différencie des concurrents.

\subsubsection{L'intemporalité}

Une bonne identité visuelle est \terme{intemporelle} : elle ne se démode pas, elle traverse les époques.

\subsubsection{La versatilité}

Une bonne identité visuelle est \terme{versatile} : elle fonctionne sur tous les supports, à toutes les tailles, en couleur et en noir et blanc.

\subsubsection{La pertinence}

Une bonne identité visuelle est \terme{pertinente} : elle correspond à la marque, à son marché, à son public.

\subsection{Les bonnes pratiques}

\begin{itemize}
    \item \terme{Définir} les composantes de la marque (nom, logo, couleurs, typographies, valeurs, personnalité, promesse).
    \item \terme{Traduire} ces composantes en identité visuelle.
    \item \terme{Vérifier} la reconnaissance.
    \item \terme{Vérifier} la cohérence.
    \item \terme{Vérifier} la distinction.
    \item \terme{Vérifier} l'intemporalité.
    \item \terme{Vérifier} la versatilité.
    \item \terme{Vérifier} la pertinence.
\end{itemize}

\subsection*{Conclusion de la section}

Une \terme{marque} est un ensemble de signes, de valeurs, de promesses et d'émotions qui permettent d'identifier un produit, un service, une organisation. Elle comprend un nom, un logo, des couleurs, des typographies, des valeurs, une personnalité, une promesse. L'\terme{identité visuelle} est la traduction graphique de la marque : logo, couleurs, typographies, éléments secondaires, règles d'application, déclinaisons. Une bonne identité visuelle est reconnaissable, cohérente, distinctive, intemporelle, versatile, pertinente.

Maintenant que nous avons défini la marque et l'identité visuelle, nous allons, dans la section suivante, distinguer \terme{logo, identité visuelle et branding}.

% ============================================================
\section{Différence entre logo, identité visuelle et branding}
% ============================================================

\subsection{Pourquoi cette distinction est-elle importante ?}

Les termes \terme{logo}, \terme{identité visuelle} et \terme{branding} sont souvent confondus. Pourtant, ils désignent des réalités différentes :

\begin{itemize}
    \item \terme{Le logo} est un élément graphique.
    \item \terme{L'identité visuelle} est un système graphique.
    \item \terme{Le branding} est une stratégie globale.
\end{itemize}

Comprendre cette distinction permet de :

\begin{itemize}
    \item \terme{Dialoguer} avec les clients.
    \item \terme{Structurer} son travail.
    \item \terme{Éviter} les malentendus.
    \item \terme{Proposer} les bons services.
    \item \terme{Facturer} correctement.
\end{itemize}

\subsection{Le logo}

\subsubsection{Définition}

Le \terme{logo} est un signe graphique qui identifie et représente une marque. Il combine généralement un symbole, une typographie, et une couleur.

\subsubsection{Rôle}

Le logo a pour rôle de :

\begin{itemize}
    \item \terme{Identifier} la marque.
    \item \terme{Différencier} la marque.
    \item \terme{Marquer} les esprits.
    \item \terme{Signaler} la présence de la marque.
\end{itemize}

\subsubsection{Limites}

Le logo seul ne suffit pas à construire une marque. Il est un point de départ, un élément parmi d'autres.

\subsection{L'identité visuelle}

\subsubsection{Définition}

L'\terme{identité visuelle} est l'ensemble des éléments graphiques qui permettent d'identifier et de représenter une marque : logo, couleurs, typographies, éléments secondaires, règles d'application.

\subsubsection{Rôle}

L'identité visuelle a pour rôle de :

\begin{itemize}
    \item \terme{Créer} un système cohérent.
    \item \terme{Assurer} la reconnaissance.
    \item \terme{Faciliter} la déclinaison.
    \item \terme{Garantir} la cohérence.
\end{itemize}

\subsubsection{Limites}

L'identité visuelle seule ne suffit pas à construire une marque. Elle est un outil, au service d'une stratégie.

\subsection{Le branding}

\subsubsection{Définition}

Le \terme{branding} (ou \emph{gestion de marque}) est l'ensemble des actions qui visent à construire, développer et gérer une marque. Il englobe :

\begin{itemize}
    \item \terme{La} stratégie de marque (positionnement, valeurs, promesse).
    \item \terme{L'}identité visuelle (logo, couleurs, typographies).
    \item \terme{La} communication (messages, canaux).
    \item \terme{L'}expérience client (service, qualité, relation).
    \item \terme{La} gestion de la marque dans le temps.
\end{itemize}

\subsubsection{Rôle}

Le branding a pour rôle de :

\begin{itemize}
    \item \terme{Définir} le positionnement.
    \item \terme{Construire} la perception.
    \item \terme{Développer} la notoriété.
    \item \terme{Créer} une relation.
    \item \terme{Fidéliser} les clients.
    \item \terme{Différencier} la marque.
\end{itemize}

\subsubsection{Le branding englobe l'identité visuelle}

Le branding \terme{englobe} l'identité visuelle : il la dépasse, il l'intègre dans une stratégie globale.

\subsection{Tableau comparatif}

\begin{table}[h!]
\centering
\begin{tabular}{@{}p{3cm}p{4cm}p{4cm}p{3.5cm}@{}}
\toprule
\textbf{Critère} & \textbf{Logo} & \textbf{Identité visuelle} & \textbf{Branding} \\
\midrule
Nature & Élément graphique & Système graphique & Stratégie globale \\
Portée & Ponctuelle & Cohérente & Globale \\
Composants & Symbole, typographie, couleur & Logo, couleurs, typographies, éléments, règles & Stratégie, identité, communication, expérience \\
Objectif & Identifier & Reconnaître & Construire une marque \\
Livrable & Fichier logo & Charte graphique & Plan de marque \\
\bottomrule
\end{tabular}
\caption{Comparaison logo, identité visuelle, branding.}
\end{table}

\subsection{Les bonnes pratiques}

\begin{itemize}
    \item \terme{Distinguer} les trois notions.
    \item \terme{Comprendre} que le logo n'est pas l'identité.
    \item \terme{Comprendre} que l'identité n'est pas le branding.
    \item \terme{Proposer} une approche globale.
    \item \terme{Dialoguer} avec les clients.
    \item \terme{Structurer} son travail.
    \item \terme{Facturer} correctement.
\end{itemize}

\subsection*{Conclusion de la section}

Le \terme{logo} est un élément graphique qui identifie une marque. L'\terme{identité visuelle} est un système graphique qui comprend le logo, les couleurs, les typographies, les éléments secondaires, les règles. Le \terme{branding} est une stratégie globale qui englobe la marque, son positionnement, sa communication, son expérience. Il est essentiel de distinguer ces trois notions, de proposer une approche globale, de facturer correctement.

Maintenant que nous avons distingué les trois notions, nous allons, dans la section suivante, apprendre à effectuer une \terme{recherche et un positionnement visuel}.

% ============================================================
\section{Recherche et positionnement visuel}
% ============================================================

\subsection{Pourquoi la recherche et le positionnement ?}

La \terme{recherche} et le \terme{positionnement visuel} sont essentiels pour :

\begin{itemize}
    \item \terme{Comprendre} le marché et la concurrence.
    \item \terme{Définir} une identité distinctive.
    \item \terme{Éviter} de ressembler aux concurrents.
    \item \terme{Aligner} l'identité avec la stratégie.
    \item \terme{Communiquer} avec le client.
    \item \terme{Construire} une marque forte.
\end{itemize}

\subsection{La recherche}

\subsubsection{La recherche sur le client}

\begin{itemize}
    \item \terme{Qui} est le client ?
    \item \terme{Que} fait-il ?
    \item \terme{Quelles} sont ses valeurs ?
    \item \terme{Quelle} est sa personnalité ?
    \item \terme{Quelle} est sa promesse ?
    \item \terme{Quel} est son public ?
\end{itemize}

\subsubsection{La recherche sur le marché}

\begin{itemize}
    \item \terme{Qui} sont les concurrents ?
    \item \terme{Quels} sont leurs logos ?
    \item \terme{Quelles} sont leurs couleurs ?
    \item \terme{Quelles} sont leurs typographies ?
    \item \terme{Quels} sont leurs positionnements ?
    \item \terme{Quelles} sont leurs forces et faiblesses ?
\end{itemize}

\subsubsection{La recherche visuelle}

\begin{itemize}
    \item \terme{Chercher} des références (Behance, Pinterest, livres).
    \item \terme{Analyser} les identités existantes.
    \item \terme{Repérer} les tendances.
    \item \terme{Identifier} les codes du secteur.
    \item \terme{Constituer} un moodboard.
\end{itemize}

\subsection{Le positionnement visuel}

\subsubsection{Définition}

Le \terme{positionnement visuel} est la place que la marque veut occuper dans l'esprit de son public, traduite en choix graphiques.

\subsubsection{Les axes de positionnement}

\begin{itemize}
    \item \terme{Classique / moderne}.
    \item \terme{Sobre / coloré}.
    \item \terme{Traditionnel / innovant}.
    \item \terme{Élégant / populaire}.
    \item \terme{Minimaliste / riche}.
    \item \terme{Sérieux / ludique}.
\end{itemize}

\subsubsection{Définir le positionnement}

\begin{enumerate}
    \item \terme{Analyser} le marché.
    \item \terme{Identifier} les positionnements existants.
    \item \terme{Identifier} les espaces libres.
    \item \terme{Définir} le positionnement souhaité.
    \item \terme{Traduire} ce positionnement en choix graphiques.
\end{enumerate}

\subsection{Les bonnes pratiques}

\begin{itemize}
    \item \terme{Rechercher} sur le client, le marché, le visuel.
    \item \terme{Analyser} les concurrents.
    \item \terme{Constituer} un moodboard.
    \item \terme{Définir} le positionnement.
    \item \terme{Traduire} en choix graphiques.
    \item \terme{Vérifier} la distinction.
    \item \terme{Vérifier} la cohérence.
    \item \terme{Communiquer} avec le client.
\end{itemize}

\subsection*{Conclusion de la section}

La \terme{recherche} et le \terme{positionnement visuel} sont essentiels pour construire une identité distinctive. La recherche porte sur le client, le marché, le visuel. Le positionnement définit la place de la marque dans l'esprit du public, traduite en choix graphiques. Il est essentiel de rechercher, d'analyser, de définir le positionnement, de vérifier la distinction.

Maintenant que nous savons effectuer une recherche et un positionnement, nous allons, dans la section suivante, apprendre à \terme{créer un logo et ses variantes}.

% ============================================================
\section{Création d'un logo et de ses variantes}
% ============================================================

\subsection{Qu'est-ce qu'un logo ?}

Nous avons déjà défini le logo (voir chapitre 20). Rappelons qu'il est l'élément graphique qui identifie et représente une marque. Il combine généralement un symbole, une typographie, et une couleur.

\subsection{Les étapes de création d'un logo}

\subsubsection{Étape 1 : Analyser le brief}

\begin{itemize}
    \item \terme{Comprendre} le client.
    \item \terme{Comprendre} le marché.
    \item \terme{Comprendre} le public.
    \item \terme{Comprendre} le positionnement.
\end{itemize}

\subsubsection{Étape 2 : Rechercher}

\begin{itemize}
    \item \terme{Chercher} des références.
    \item \terme{Analyser} les concurrents.
    \item \terme{Constituer} un moodboard.
\end{itemize}

\subsubsection{Étape 3 : Esquisser}

\begin{itemize}
    \item \terme{Dessiner} plusieurs pistes.
    \item \terme{Explorer} différentes directions.
    \item \terme{Sélectionner} les meilleures.
\end{itemize}

\subsubsection{Étape 4 : Créer}

\begin{itemize}
    \item \terme{Créer} le logo dans Illustrator.
    \item \terme{Ajuster} les formes.
    \item \terme{Ajuster} les couleurs.
    \item \terme{Ajuster} la typographie.
\end{itemize}

\subsubsection{Étape 5 : Vérifier}

\begin{itemize}
    \item \terme{Vérifier} la lisibilité.
    \item \terme{Vérifier} l'équilibre.
    \item \terme{Vérifier} la cohérence.
    \item \terme{Vérifier} la pertinence.
\end{itemize}

\subsubsection{Étape 6 : Décliner}

\begin{itemize}
    \item \terme{Créer} les variantes.
    \item \terme{Créer} les versions monochromes.
    \item \terme{Créer} les versions adaptées.
\end{itemize}

\subsection{Les variantes du logo}

\subsubsection{Version horizontale}

La version \terme{horizontale} présente le symbole et le texte côte à côte. Elle est utilisée dans les espaces larges (en-tête de site, bannière, signature email).

\subsubsection{Version verticale}

La version \terme{verticale} présente le symbole au-dessus du texte. Elle est utilisée dans les espaces étroits (avatar, favicon, carte de visite carrée).

\subsubsection{Version carrée}

La version \terme{carrée} présente le logo dans un carré. Elle est utilisée pour les avatars, les favicons, les applications.

\subsubsection{Symbole seul}

Le \terme{symbole seul} est utilisé pour les petites tailles (favicon, avatar, icône). Il doit être reconnaissable.

\subsubsection{Texte seul}

Le \terme{texte seul} est utilisé pour les contextes où le symbole n'est pas nécessaire (signature, mention). Il doit être lisible.

\subsubsection{Version monochrome}

La version \terme{monochrome} est utilisée pour les impressions en noir et blanc, les gravures, les fax.

\subsection{Les règles du logo}

\subsubsection{La zone de protection}

La \terme{zone de protection} est l'espace minimal autour du logo, où aucun autre élément ne doit être placé. Elle garantit la lisibilité et l'impact du logo.

\subsubsection{La taille minimale}

La \terme{taille minimale} est la plus petite taille à laquelle le logo reste lisible. Elle doit être respectée.

\subsubsection{Les usages}

Les \terme{usages} définissent quand et comment utiliser chaque variante du logo. Ils sont documentés dans la charte graphique.

\subsection{Les bonnes pratiques}

\begin{itemize}
    \item \terme{Analyser} le brief.
    \item \terme{Rechercher} des références.
    \item \terme{Esquisser} plusieurs pistes.
    \item \terme{Créer} le logo.
    \item \terme{Vérifier} la qualité.
    \item \terme{Décliner} en variantes.
    \item \terme{Définir} la zone de protection.
    \item \terme{Définir} la taille minimale.
    \item \terme{Documenter} les usages.
\end{itemize}

\subsection*{Conclusion de la section}

Le \terme{logo} est l'élément graphique qui identifie une marque. Sa création suit plusieurs étapes : analyser le brief, rechercher, esquisser, créer, vérifier, décliner. Il existe plusieurs variantes : horizontale, verticale, carrée, symbole seul, texte seul, monochrome. Il est essentiel de définir la zone de protection, la taille minimale, les usages.

Maintenant que nous savons créer un logo et ses variantes, nous allons, dans la section suivante, apprendre à \terme{choisir les couleurs de marque}.

% ============================================================
\section{Choix des couleurs de marque}
% ============================================================

\subsection{Pourquoi les couleurs sont-elles essentielles ?}

Les \terme{couleurs} sont un élément fort de l'identité visuelle. Elles permettent de :

\begin{itemize}
    \item \terme{Attirer} l'attention.
    \item \terme{Transmettre} une émotion.
    \item \terme{Signifier} des valeurs.
    \item \terme{Différencier} la marque.
    \item \terme{Créer} une reconnaissance.
    \item \terme{Unifier} les supports.
\end{itemize}

\subsection{Les principes de choix des couleurs}

\subsubsection{La cohérence}

Les couleurs doivent être \terme{cohérentes} avec les valeurs, le positionnement, le public de la marque.

\subsubsection{La distinction}

Les couleurs doivent être \terme{distinctives} : différentes des concurrents.

\subsubsection{La mémorisation}

Les couleurs doivent être \terme{mémorables} : associées à la marque.

\subsubsection{La pertinence}

Les couleurs doivent être \terme{pertinentes} : adaptées au secteur, à la culture, au public.

\subsubsection{La versatilité}

Les couleurs doivent être \terme{versatiles} : fonctionner sur tous les supports (écran, impression, textile).

\subsection{La palette de marque}

\subsubsection{La couleur principale}

La \terme{couleur principale} est la couleur dominante. Elle définit l'ambiance. Elle est souvent utilisée pour le logo, les éléments clés.

\subsubsection{Les couleurs secondaires}

Les \terme{couleurs secondaires} complètent la couleur principale. Elles sont utilisées pour les éléments importants, mais pas dominants.

\subsubsection{Les couleurs d'accent}

Les \terme{couleurs d'accent} sont utilisées pour attirer l'attention sur des éléments spécifiques (boutons, liens, titres).

\subsubsection{Les couleurs neutres}

Les \terme{couleurs neutres} (blanc, noir, gris) servent de base, de fond, de texte.

\subsection{Choisir les couleurs de marque}

\subsubsection{Étape 1 : Analyser}

\begin{itemize}
    \item \terme{Quelles} sont les valeurs de la marque ?
    \item \terme{Quelle} est la personnalité ?
    \item \terme{Quel} est le positionnement ?
    \item \terme{Quel} est le public ?
    \item \terme{Quels} sont les concurrents ?
\end{itemize}

\subsubsection{Étape 2 : Choisir la couleur principale}

\begin{itemize}
    \item \terme{Choisir} une couleur qui correspond à l'identité.
    \item \terme{Vérifier} les significations culturelles.
    \item \terme{Vérifier} la différenciation.
\end{itemize}

\subsubsection{Étape 3 : Choisir les couleurs secondaires}

\begin{itemize}
    \item \terme{Choisir} des couleurs qui complètent la couleur principale.
    \item \terme{Utiliser} une harmonie (analogue, complémentaire, triadique).
    \item \terme{Vérifier} la cohérence.
\end{itemize}

\subsubsection{Étape 4 : Choisir les couleurs d'accent}

\begin{itemize}
    \item \terme{Choisir} une couleur vive, contrastée.
    \item \terme{Utiliser} avec parcimonie.
    \item \terme{Vérifier} la lisibilité.
\end{itemize}

\subsubsection{Étape 5 : Choisir les couleurs neutres}

\begin{itemize}
    \item \terme{Choisir} du blanc, du noir, du gris.
    \item \terme{Définir} les usages (fond, texte, bordures).
\end{itemize}

\subsubsection{Étape 6 : Définir les codes}

\begin{itemize}
    \item \terme{Définir} les codes RVB (pour l'écran).
    \item \terme{Définir} les codes CMJN (pour l'impression).
    \item \terme{Définir} les codes Pantone (pour l'impression).
    \item \terme{Définir} les codes HEX (pour le web).
\end{itemize}

\subsubsection{Étape 7 : Tester}

\begin{itemize}
    \item \terme{Tester} la palette sur différents supports.
    \item \terme{Vérifier} la cohérence.
    \item \terme{Vérifier} la lisibilité.
    \item \terme{Vérifier} le contraste.
\end{itemize}

\subsection{Les bonnes pratiques}

\begin{itemize}
    \item \terme{Analyser} la marque.
    \item \terme{Choisir} une couleur principale.
    \item \terme{Choisir} des couleurs secondaires.
    \item \terme{Choisir} des accents.
    \item \terme{Définir} les codes.
    \item \terme{Tester} sur différents supports.
    \item \terme{Vérifier} la cohérence.
    \item \terme{Documenter} dans la charte.
\end{itemize}

\subsection*{Conclusion de la section}

Les \terme{couleurs} sont un élément fort de l'identité visuelle. Elles doivent être cohérentes, distinctives, mémorables, pertinentes, versatiles. La palette comprend une couleur principale, des couleurs secondaires, des accents, des neutres. Le choix suit plusieurs étapes : analyser, choisir la couleur principale, les secondaires, les accents, les neutres, définir les codes, tester. Il est essentiel de choisir une palette cohérente, de définir les codes, de tester sur différents supports.

Maintenant que nous savons choisir les couleurs, nous allons, dans la section suivante, apprendre à \terme{définir le système typographique}.

% ============================================================
\section{Définition du système typographique}
% ============================================================

\subsection{Pourquoi le système typographique ?}

Le \terme{système typographique} est essentiel pour :

\begin{itemize}
    \item \terme{Transmettre} la personnalité de la marque.
    \item \terme{Assurer} la lisibilité.
    \item \terme{Créer} une cohérence.
    \item \terme{Faciliter} la déclinaison.
    \item \terme{Différencier} la marque.
    \item \terme{Unifier} les supports.
\end{itemize}

\subsection{Les composants du système typographique}

\subsubsection{La police de titre}

La \terme{police de titre} est utilisée pour les titres, les accroches, les slogans. Elle doit être :

\begin{itemize}
    \item \terme{Distinctive} : reconnaissable.
    \item \terme{Impactante} : attire l'attention.
    \item \terme{Lisible} : même en grand.
    \item \terme{Cohérente} : avec la personnalité.
\end{itemize}

\subsubsection{La police de texte}

La \terme{police de texte} est utilisée pour le texte courant, les paragraphes, les descriptions. Elle doit être :

\begin{itemize}
    \item \terme{Lisible} : à toutes les tailles.
    \item \terme{Neutre} : ne pas distraire.
    \item \terme{Cohérente} : avec la police de titre.
    \item \terme{Versatile} : sur tous les supports.
\end{itemize}

\subsubsection{La police d'accent}

La \terme{police d'accent} est utilisée pour les citations, les légendes, les mentions. Elle peut être :

\begin{itemize}
    \item \terme{Une} variante de la police de titre.
    \item \terme{Une} police complémentaire.
    \item \terme{Une} police décorative.
\end{itemize}

\subsection{Choisir les typographies}

\subsubsection{Étape 1 : Analyser}

\begin{itemize}
    \item \terme{Quelle} est la personnalité de la marque ?
    \item \terme{Quel} est le positionnement ?
    \item \terme{Quel} est le public ?
\end{itemize}

\subsubsection{Étape 2 : Choisir la police de titre}

\begin{itemize}
    \item \terme{Choisir} une police distinctive.
    \item \terme{Vérifier} la lisibilité.
    \item \terme{Vérifier} la cohérence.
\end{itemize}

\subsubsection{Étape 3 : Choisir la police de texte}

\begin{itemize}
    \item \terme{Choisir} une police lisible.
    \item \terme{Vérifier} la complémentarité.
    \item \terme{Vérifier} la versatilité.
\end{itemize}

\subsubsection{Étape 4 : Définir les usages}

\begin{itemize}
    \item \terme{Définir} les tailles (titres, sous-titres, texte).
    \item \terme{Définir} les graisses (normal, gras).
    \item \terme{Définir} les interlignages.
    \item \terme{Définir} les alignements.
\end{itemize}

\subsubsection{Étape 5 : Tester}

\begin{itemize}
    \item \terme{Tester} sur différents supports.
    \item \terme{Vérifier} la lisibilité.
    \item \terme{Vérifier} la cohérence.
\end{itemize}

\subsection{Les bonnes pratiques}

\begin{itemize}
    \item \terme{Analyser} la marque.
    \item \terme{Choisir} une police de titre.
    \item \terme{Choisir} une police de texte.
    \item \terme{Définir} les usages.
    \item \terme{Tester} sur différents supports.
    \item \terme{Vérifier} la lisibilité.
    \item \terme{Vérifier} la cohérence.
    \item \terme{Documenter} dans la charte.
\end{itemize}

\subsection*{Conclusion de la section}

Le \terme{système typographique} comprend une police de titre, une police de texte, une police d'accent. Il doit transmettre la personnalité, assurer la lisibilité, créer une cohérence. Le choix suit plusieurs étapes : analyser, choisir les polices, définir les usages, tester. Il est essentiel de choisir des polices cohérentes, de définir les usages, de tester sur différents supports.

Maintenant que nous savons définir le système typographique, nous allons, dans la section suivante, apprendre à \terme{créer des éléments graphiques secondaires}.

% ============================================================
\section{Création d'éléments graphiques secondaires}
% ============================================================

\subsection{Qu'est-ce qu'un élément graphique secondaire ?}

\begin{quote}
\textbf{Un élément graphique secondaire est un élément visuel qui complète le logo, les couleurs et les typographies : icônes, pictogrammes, motifs, illustrations, formes, textures.}
\end{quote}

Les éléments secondaires enrichissent l'identité visuelle, lui donnent de la personnalité, de la cohérence, de la distinction.

\subsection{Les types d'éléments secondaires}

\subsubsection{Les icônes}

Les \terme{icônes} sont des représentations graphiques simplifiées. Elles sont utilisées dans les interfaces, la signalétique, les infographies.

\subsubsection{Les pictogrammes}

Les \terme{pictogrammes} sont des représentations graphiques universelles. Ils sont utilisés dans la signalétique, les modes d'emploi, les infographies.

\subsubsection{Les motifs}

Les \terme{motifs} sont des formes répétées. Ils sont utilisés pour les fonds, les packagings, les papeteries.

\subsubsection{Les illustrations}

Les \terme{illustrations} sont des images dessinées. Elles sont utilisées pour raconter, expliquer, décorer.

\subsubsection{Les formes graphiques}

Les \terme{formes graphiques} (lignes, cercles, triangles) enrichissent les compositions.

\subsubsection{Les textures}

Les \terme{textures} (papier, tissu, bois) ajoutent de la matière.

\subsection{Créer des éléments secondaires}

\subsubsection{Étape 1 : Analyser la marque}

\begin{itemize}
    \item \terme{Quelles} sont les valeurs ?
    \item \terme{Quelle} est la personnalité ?
    \item \terme{Quel} est le positionnement ?
\end{itemize}

\subsubsection{Étape 2 : Choisir le type d'élément}

\begin{itemize}
    \item \terme{Choisir} le type (icônes, pictogrammes, motifs, illustrations).
    \item \terme{Choisir} le style (flat, ligne, plein).
\end{itemize}

\subsubsection{Étape 3 : Créer}

\begin{itemize}
    \item \terme{Dessiner} les éléments dans Illustrator.
    \item \terme{Utiliser} les formes de base.
    \item \terme{Simplifier} les tracés.
\end{itemize}

\subsubsection{Étape 4 : Vérifier}

\begin{itemize}
    \item \terme{Vérifier} la cohérence avec le logo.
    \item \terme{Vérifier} la cohérence avec les couleurs.
    \item \terme{Vérifier} la cohérence avec les typographies.
\end{itemize}

\subsubsection{Étape 5 : Documenter}

\begin{itemize}
    \item \terme{Documenter} les usages.
    \item \terme{Documenter} les règles.
    \item \terme{Documenter} les exemples.
\end{itemize}

\subsection{Les bonnes pratiques}

\begin{itemize}
    \item \terme{Analyser} la marque.
    \item \terme{Choisir} le type d'élément.
    \item \terme{Créer} les éléments.
    \item \terme{Vérifier} la cohérence.
    \item \terme{Documenter} les usages.
    \item \terme{Maintenir} un style cohérent.
    \item \terme{Éviter} la surcharge.
    \item \terme{Tester} sur différents supports.
\end{itemize}

\subsection*{Conclusion de la section}

Les \terme{éléments graphiques secondaires} complètent le logo, les couleurs, les typographies : icônes, pictogrammes, motifs, illustrations, formes, textures. Ils enrichissent l'identité visuelle, lui donnent de la personnalité, de la cohérence, de la distinction. La création suit plusieurs étapes : analyser, choisir, créer, vérifier, documenter. Il est essentiel de maintenir un style cohérent, d'éviter la surcharge, de tester sur différents supports.

Maintenant que nous savons créer des éléments secondaires, nous allons, dans la section suivante, apprendre à \terme{appliquer l'identité sur les supports de communication}.

% ============================================================
\section{Applications sur les supports de communication}
% ============================================================

\subsection{Pourquoi appliquer l'identité sur les supports ?}

L'\terme{application} de l'identité sur les supports est essentielle pour :

\begin{itemize}
    \item \terme{Assurer} la cohérence.
    \item \terme{Faciliter} la reconnaissance.
    \item \terme{Démontrer} la versatilité.
    \item \terme{Fournir} des exemples au client.
    \item \terme{Documenter} la charte.
    \item \terme{Préparer} la production.
\end{itemize}

\subsection{Les supports d'application}

\subsubsection{La papeterie}

\begin{itemize}
    \item \terme{Carte} de visite.
    \item \terme{Papier} à en-tête.
    \item \terme{Enveloppe}.
    \item \terme{Pochette}.
    \item \terme{Bloc-notes}.
\end{itemize}

\subsubsection{Les supports numériques}

\begin{itemize}
    \item \terme{Site} web.
    \item \terme{Application} mobile.
    \item \terme{Réseaux} sociaux.
    \item \terme{Newsletter}.
    \item \terme{Bannière} web.
\end{itemize}

\subsubsection{Les supports publicitaires}

\begin{itemize}
    \item \terme{Affiche}.
    \item \terme{Flyer}.
    \item \terme{Dépliant}.
    \item \terme{Publicité} en ligne.
\end{itemize}

\subsubsection{Le packaging}

\begin{itemize}
    \item \terme{Boîte}.
    \item \terme{Étiquette}.
    \item \terme{Sachet}.
    \item \terme{Flacon}.
\end{itemize}

\subsubsection{La signalétique}

\begin{itemize}
    \item \terme{Panneau}.
    \item \terme{Enseigne}.
    \item \terme{Signalétique} intérieure.
    \item \terme{Signalétique} extérieure.
\end{itemize}

\subsubsection{Les objets promotionnels}

\begin{itemize}
    \item \terme{T-shirt}.
    \item \terme{Mug}.
    \item \terme{Stylo}.
    \item \terme{Sac}.
\end{itemize}

\subsection{Appliquer l'identité}

\subsubsection{Étape 1 : Choisir les supports}

\begin{itemize}
    \item \terme{Choisir} les supports pertinents.
    \item \terme{Hiérarchiser} les supports.
    \item \terme{Préparer} les déclinaisons.
\end{itemize}

\subsubsection{Étape 2 : Décliner}

\begin{itemize}
    \item \terme{Décliner} le logo.
    \item \terme{Décliner} les couleurs.
    \item \terme{Décliner} les typographies.
    \item \terme{Décliner} les éléments secondaires.
\end{itemize}

\subsubsection{Étape 3 : Vérifier}

\begin{itemize}
    \item \terme{Vérifier} la cohérence.
    \item \terme{Vérifier} la lisibilité.
    \item \terme{Vérifier} la qualité.
\end{itemize}

\subsubsection{Étape 4 : Documenter}

\begin{itemize}
    \item \terme{Documenter} les applications.
    \item \terme{Documenter} les règles.
    \item \terme{Documenter} les exemples.
\end{itemize}

\subsection{Les bonnes pratiques}

\begin{itemize}
    \item \terme{Choisir} les supports pertinents.
    \item \terme{Décliner} avec soin.
    \item \terme{Respecter} la charte.
    \item \terme{Vérifier} la cohérence.
    \item \terme{Vérifier} la lisibilité.
    \item \terme{Documenter} les applications.
    \item \terme{Tester} sur différents supports.
    \item \terme{Préparer} la production.
\end{itemize}

\subsection*{Conclusion de la section}

L'\terme{application} de l'identité sur les supports est essentielle. Les supports sont : papeterie, numérique, publicitaire, packaging, signalétique, objets promotionnels. L'application suit plusieurs étapes : choisir les supports, décliner, vérifier, documenter. Il est essentiel de choisir les supports pertinents, de décliner avec soin, de vérifier la cohérence.

Maintenant que nous savons appliquer l'identité, nous allons, dans la section suivante, apprendre à \terme{élaborer une charte graphique}.

% ============================================================
\section{Élaboration d'une charte graphique}
% ============================================================

\subsection{Qu'est-ce qu'une charte graphique ?}

\begin{quote}
\textbf{Une charte graphique est un document qui définit les règles d'utilisation de l'identité visuelle d'une marque : logo, couleurs, typographies, éléments secondaires, applications.}
\end{quote}

La charte graphique est essentielle pour :

\begin{itemize}
    \item \terme{Assurer} la cohérence.
    \item \terme{Faciliter} la production.
    \item \terme{Former} les équipes.
    \item \terme{Contrôler} la qualité.
    \item \terme{Protéger} la marque.
    \item \terme{Transmettre} l'identité.
\end{itemize}

\subsection{Le contenu d'une charte graphique}

\subsubsection{La présentation de la marque}

\begin{itemize}
    \item \terme{Histoire} de la marque.
    \item \terme{Valeurs}.
    \item \terme{Personnalité}.
    \item \terme{Positionnement}.
    \item \terme{Public}.
\end{itemize}

\subsubsection{Le logo}

\begin{itemize}
    \item \terme{Présentation} du logo.
    \item \terme{Variantes} (horizontale, verticale, carrée, symbole seul, texte seul).
    \item \terme{Versions} monochromes.
    \item \terme{Zone} de protection.
    \item \terme{Taille} minimale.
    \item \terme{Usages}.
    \item \terme{Interdits} (déformations, couleurs non autorisées).
\end{itemize}

\subsubsection{Les couleurs}

\begin{itemize}
    \item \terme{Palette} (principale, secondaires, accents, neutres).
    \item \terme{Codes} (RVB, CMJN, Pantone, HEX).
    \item \terme{Usages}.
    \item \terme{Combinaisons} autorisées.
\end{itemize}

\subsubsection{Les typographies}

\begin{itemize}
    \item \terme{Polices} (titre, texte, accent).
    \item \terme{Tailles}.
    \item \terme{Graisses}.
    \item \terme{Interlignages}.
    \item \terme{Usages}.
    \item \terme{Exemples}.
\end{itemize}

\subsubsection{Les éléments secondaires}

\begin{itemize}
    \item \terme{Icônes}.
    \item \terme{Pictogrammes}.
    \item \terme{Motifs}.
    \item \terme{Illustrations}.
    \item \terme{Usages}.
\end{itemize}

\subsubsection{Les applications}

\begin{itemize}
    \item \terme{Papeterie}.
    \item \terme{Web}.
    \item \terme{Publicité}.
    \item \terme{Packaging}.
    \item \terme{Signalétique}.
    \item \terme{Exemples}.
\end{itemize}

\subsubsection{Les règles}

\begin{itemize}
    \item \terme{Règles} d'utilisation.
    \item \terme{Interdits}.
    \item \terme{Bonnes} pratiques.
\end{itemize}

\subsection{Élaborer une charte graphique}

\subsubsection{Étape 1 : Rassembler les éléments}

\begin{itemize}
    \item \terme{Rassembler} le logo, les variantes.
    \item \terme{Rassembler} la palette.
    \item \terme{Rassembler} les typographies.
    \item \terme{Rassembler} les éléments secondaires.
    \item \terme{Rassembler} les applications.
\end{itemize}

\subsubsection{Étape 2 : Structurer}

\begin{itemize}
    \item \terme{Structurer} les sections.
    \item \terme{Définir} l'ordre.
    \item \terme{Créer} une cohérence visuelle.
\end{itemize}

\subsubsection{Étape 3 : Rédiger}

\begin{itemize}
    \item \terme{Expliquer} les choix.
    \item \terme{Définir} les usages.
    \item \terme{Donner} des exemples.
    \item \terme{Préciser} les interdits.
\end{itemize}

\subsubsection{Étape 4 : Illustrer}

\begin{itemize}
    \item \terme{Présenter} les visuels.
    \item \terme{Montrer} les applications.
    \item \terme{Utiliser} des mockups.
\end{itemize}

\subsubsection{Étape 5 : Mettre en page}

\begin{itemize}
    \item \terme{Choisir} un format (A4, paysage, portrait).
    \item \terme{Choisir} une police.
    \item \terme{Créer} une mise en page claire.
\end{itemize}

\subsubsection{Étape 6 : Valider}

\begin{itemize}
    \item \terme{Valider} avec le client.
    \item \terme{Corriger} si nécessaire.
    \item \terme{Finaliser}.
\end{itemize}

\subsection{Les bonnes pratiques}

\begin{itemize}
    \item \terme{Rassembler} les éléments.
    \item \terme{Structurer} la charte.
    \item \terme{Rédiger} clairement.
    \item \terme{Illustrer} avec des exemples.
    \item \terme{Mettre} en page.
    \item \terme{Valider} avec le client.
    \item \terme{Diffuser} la charte.
    \item \terme{Former} les équipes.
\end{itemize}

\subsection*{Conclusion de la section}

Une \terme{charte graphique} est un document qui définit les règles d'utilisation de l'identité visuelle. Elle comprend la présentation de la marque, le logo, les couleurs, les typographies, les éléments secondaires, les applications, les règles. L'élaboration suit plusieurs étapes : rassembler, structurer, rédiger, illustrer, mettre en page, valider. Il est essentiel de rassembler, de structurer, de rédiger, d'illustrer, de valider.

Maintenant que nous savons élaborer une charte graphique, nous allons, dans la dernière section de ce chapitre, apprendre à \terme{présenter un projet d'identité visuelle}.

% ============================================================
\section{Présentation d'un projet d'identité visuelle}
% ============================================================

\subsection{Pourquoi présenter un projet d'identité visuelle ?}

La \terme{présentation} d'un projet d'identité visuelle est essentielle pour :

\begin{itemize}
    \item \terme{Communiquer} la démarche.
    \item \terme{Expliquer} les choix.
    \item \terme{Convaincre} le client.
    \item \terme{Démontrer} la cohérence.
    \item \terme{Faciliter} la validation.
    \item \terme{Documenter} le projet.
\end{itemize}

Une identité visuelle bien présentée est une identité visuelle comprise, acceptée, valorisée.

\subsection{Les éléments de la présentation}

\subsubsection{La couverture}

\begin{itemize}
    \item \terme{Nom} de la marque.
    \item \terme{Titre} du projet.
    \item \terme{Logo}.
    \item \terme{Date}.
\end{itemize}

\subsubsection{Le brief}

\begin{itemize}
    \item \terme{Objectifs}.
    \item \terme{Public}.
    \item \terme{Contraintes}.
\end{itemize}

\subsubsection{La démarche}

\begin{itemize}
    \item \terme{Recherche}.
    \item \terme{Inspirations}.
    \item \terme{Esquisses}.
    \item \terme{Concepts}.
\end{itemize}

\subsubsection{Le logo}

\begin{itemize}
    \item \terme{Présentation}.
    \item \terme{Variantes}.
    \item \terme{Versions} monochromes.
    \item \terme{Zones} de protection.
    \item \terme{Usages}.
\end{itemize}

\subsubsection{Les couleurs}

\begin{itemize}
    \item \terme{Palette}.
    \item \terme{Codes}.
    \item \terme{Usages}.
\end{itemize}

\subsubsection{Les typographies}

\begin{itemize}
    \item \terme{Polices}.
    \item \terme{Usages}.
    \item \terme{Exemples}.
\end{itemize}

\subsubsection{Les éléments secondaires}

\begin{itemize}
    \item \terme{Icônes}.
    \item \terme{Pictogrammes}.
    \item \terme{Motifs}.
\end{itemize}

\subsubsection{Les applications}

\begin{itemize}
    \item \terme{Papeterie}.
    \item \terme{Web}.
    \item \terme{Publicité}.
    \item \terme{Packaging}.
    \item \terme{Signalétique}.
\end{itemize}

\subsubsection{La charte}

\begin{itemize}
    \item \terme{Présentation} de la charte.
    \item \terme{Règles}.
    \item \terme{Exemples}.
\end{itemize}

\subsection{Présenter}

\subsubsection{Étape 1 : Structurer}

\begin{itemize}
    \item \terme{Organiser} la présentation.
    \item \terme{Définir} les sections.
    \item \terme{Créer} une cohérence.
\end{itemize}

\subsubsection{Étape 2 : Rédiger}

\begin{itemize}
    \item \terme{Expliquer} les choix.
    \item \terme{Justifier} les décisions.
    \item \terme{Démontrer} la cohérence.
\end{itemize}

\subsubsection{Étape 3 : Illustrer}

\begin{itemize}
    \item \terme{Présenter} les visuels.
    \item \terme{Montrer} les applications.
    \item \terme{Utiliser} des mockups.
\end{itemize}

\subsubsection{Étape 4 : Présenter oralement}

\begin{itemize}
    \item \terme{Présenter} oralement.
    \item \terme{Expliquer} la démarche.
    \item \terme{Répondre} aux questions.
\end{itemize}

\subsubsection{Étape 5 : Livrer}

\begin{itemize}
    \item \terme{Livrer} les fichiers.
    \item \terme{Livrer} la charte.
    \item \terme{Livrer} les exports.
\end{itemize}

\subsection{Les bonnes pratiques}

\begin{itemize}
    \item \terme{Structurer} la présentation.
    \item \terme{Expliquer} les choix.
    \item \terme{Justifier} les décisions.
    \item \terme{Illustrer} avec des visuels.
    \item \terme{Présenter} les applications.
    \item \terme{Utiliser} des mockups.
    \item \terme{Présenter} oralement.
    \item \terme{Livrer} les fichiers.
\end{itemize}

\subsection*{Conclusion de la section}

La \terme{présentation} d'un projet d'identité visuelle est essentielle pour communiquer la démarche, expliquer les choix, convaincre le client. Elle comprend la couverture, le brief, la démarche, le logo, les couleurs, les typographies, les éléments secondaires, les applications, la charte. La présentation suit plusieurs étapes : structurer, rédiger, illustrer, présenter oralement, livrer. Il est essentiel de structurer, d'expliquer, d'illustrer, de présenter oralement.

% ============================================================
\section*{Conclusion du chapitre}
\addcontentsline{toc}{section}{Conclusion du chapitre}
% ============================================================

Nous voici au terme du chapitre 25. Nous avons appris à maîtriser le branding et les systèmes d'identité visuelle :

\begin{itemize}
    \item \terme{Définition de la marque et de l'identité visuelle} : composantes (nom, logo, couleurs, typographies, valeurs, personnalité, promesse), qualités (reconnaissance, cohérence, distinction, intemporalité, versatilité, pertinence).
    \item \terme{Différence entre logo, identité visuelle et branding} : le logo est un élément, l'identité un système, le branding une stratégie.
    \item \terme{Recherche et positionnement visuel} : recherche sur le client, le marché, le visuel ; axes de positionnement.
    \item \terme{Création d'un logo et de ses variantes} : étapes, variantes, zone de protection, taille minimale.
    \item \terme{Choix des couleurs de marque} : principes, palette, codes.
    \item \terme{Définition du système typographique} : police de titre, de texte, d'accent ; usages.
    \item \terme{Création d'éléments graphiques secondaires} : icônes, pictogrammes, motifs, illustrations.
    \item \terme{Applications sur les supports} : papeterie, numérique, publicité, packaging, signalétique.
    \item \terme{Élaboration d'une charte graphique} : contenu, étapes.
    \item \terme{Présentation d'un projet d'identité visuelle} : éléments, étapes.
\end{itemize}

Le branding est bien plus qu'un logo, plus qu'une palette de couleurs, plus qu'une typographie : c'est l'ensemble des perceptions, des émotions et des associations qu'une marque suscite chez son public. Le branding ne se limite pas au design : il englobe la stratégie, le positionnement, la communication, l'expérience client. Mais le design graphique en est l'un des piliers : c'est lui qui donne une forme visible, tangible, mémorable à la marque.

Un système d'identité visuelle est l'ensemble cohérent des éléments graphiques qui permettent de reconnaître une marque : logo, couleurs, typographies, éléments secondaires, règles d'application. Ce système doit être pensé, construit, documenté, pour être utilisé de manière cohérente sur tous les supports. Une identité visuelle réussie est reconnaissable, cohérente, intemporelle, adaptable. Une identité visuelle ratée est incohérente, générique, datée, difficile à décliner.

Ce chapitre t'a donné les bases du branding et des systèmes d'identité visuelle. Mais la maîtrise vient avec la pratique. N'hésite pas à t'entraîner sur des projets variés : marques fictives, refontes, identités complètes. Chaque projet est une occasion de progresser.

Dans le chapitre suivant, nous aborderons le \terme{photomontage publicitaire avancé}. Tu apprendras à analyser un concept publicitaire, à rechercher et sélectionner des photographies, à construire une composition, à réaliser un détourage avancé, à harmoniser la lumière, à gérer les ombres et les reflets, à intégrer un produit, à traiter la colorimétrie, à intégrer la typographie, à finaliser et à contrôler la qualité. Ces compétences te permettront de créer des visuels publicitaires spectaculaires.

\vspace{2em}

\begin{center}
\textit{Fin du chapitre 25 --- Branding et systèmes d'identité visuelle.}

\vspace{1em}

\textbf{Prochain chapitre : Chapitre 26 --- Photomontage publicitaire avancé.}
\end{center}